\par\Needspace{6\baselineskip}\medskip
\hypertarget{native-127-article094}{}\pdfbookmark[2]{94.}{article-094}
\begin{center}94.\end{center}
The treatment of equations 12 and $12^{\ast}$ explained in Arts. 91, 92, and 93 rests on the assumption that the angle $g$ is not too large and is certainly below the limit $66^0\,25'$, beyond which we do not extend Table III. When this assumption does not hold, those equations need no such elaborate treatment. They can always be solved \textit{without changing their form}, very reliably and conveniently, by trial. \textit{Reliably}, because the value of the expression $\dfrac{2g-\sin 2g}{\sin g^3}\mathpunct{\raisebox{.45ex}{,}}$ in which $2g$ is of course to be expressed in parts of the radius, can be computed from trigonometric tables \textit{with full precision} for larger values of $g$, which is impossible while $g$ is a small angle. \textit{Conveniently}, because heliocentric positions separated by so large an interval will scarcely ever be used to determine an orbit that is still entirely unknown, and from any knowledge of the orbit an approximate value of $g$ follows almost effortlessly by equation 1 or 3 of Art. 88. Finally, from the approximate value of $g$ a few trials will always yield the corrected value, satisfying equation 12 or $12^{\ast}$ to any desired precision. Moreover, whenever the two given heliocentric positions span more than one complete revolution, we must bear in mind that the eccentric anomaly will have completed just as many full revolutions, so that the angles $E'-E$ and $v'-v$ both \SourceLeaf{128}lie either between 0 and $360^0$ or between corresponding multiples of a full revolution. Thus $f$ and $g$ both lie either between 0 and $180^0$ or between corresponding multiples of a half revolution. Finally, if the orbit were entirely unknown, and it were therefore uncertain whether the celestial body, in passing from the first radius vector to the second, had traversed only part of a revolution or had in addition completed one or more full revolutions, our problem might in some cases admit several different solutions. We pass over this case here, as it will scarcely ever occur in practice.
